% ===========================================================================
% Chapitre 2 — Ingénierie des exigences.
%
% Les identifiants FR-* sont ceux du projet. Ils sont repris tels quels au
% chapitre 4 (traçabilité d'implémentation) et au chapitre 5 (validation) :
% ne les renuméroter qu'en cascade.
% ===========================================================================
\chapter{Ingénierie des exigences}
\label{chap:requirements}

\section{Introduction}
\label{sec:req_intro}

Ce chapitre transforme le cadrage du chapitre précédent en exigences
vérifiables. Il identifie les acteurs et leurs interactions, expose la méthode
de recueil, énonce les exigences fonctionnelles par domaine et les exigences
non fonctionnelles avec leurs critères d'acceptation, décrit les cas
d'utilisation principaux, et établit le modèle de traçabilité qui relie chaque
exigence à sa réalisation et à sa validation.

\section{Parties prenantes et acteurs}
\label{sec:actors}

Le tableau~\ref{tab:actors} recense les acteurs de la plateforme et leurs
interactions principales.

\begin{table}[htbp]
  \centering
  \caption{Acteurs de PIQBIT et interactions principales}
  \label{tab:actors}
  \begin{tabularx}{\textwidth}{L{2.8cm} L{4.0cm} X}
    \toprule
    \textbf{Acteur} & \textbf{Description} & \textbf{Interactions principales} \\
    \midrule
    Visiteur
      & Personne non authentifiée parcourant le site.
      & Consulter les offres publiées, rechercher, créer un compte. \\
    Candidat
      & Utilisateur authentifié déposant des candidatures.
      & Déposer un CV, postuler, suivre ses candidatures, consulter l'analyse de son CV, retirer une candidature. \\
    Porteur de lien
      & Personne agissant sans compte, au moyen d'un jeton reçu par courriel.
      & Confirmer un créneau d'entretien, passer le test technique, signer ou refuser un contrat. \\
    Recruteur
      & Rôle interne conduisant l'évaluation des candidats.
      & Consulter le classement, planifier et évaluer des entretiens, consulter les analyses de CV. \\
    Responsable RH
      & Rôle interne pilotant le recrutement et l'embauche.
      & Publier des offres, décider, éditer des contrats, piloter employés, budgets et carrières. \\
    Administrateur
      & Rôle interne d'administration du système.
      & Gérer les comptes et les rôles, consulter le journal d'audit. \\
    Traitement de fond
      & Processus asynchrones de la plateforme.
      & Analyser les CV, indexer les vecteurs, expirer les invitations, envoyer les rappels. \\
    Service externe
      & Dépendance tierce, toujours facultative.
      & Messagerie, agenda, visioconférence, modèles de langage, dépôt de modèles. \\
    \bottomrule
  \end{tabularx}
\end{table}

Le modèle de rôles implémenté comporte quatre rôles — \texttt{ADMIN},
\texttt{RH}, \texttt{RECRUTEUR}, \texttt{CANDIDAT} — liés par un héritage de
permissions : le responsable RH hérite des droits du recruteur, et
l'administrateur de ceux du responsable RH. Le « visiteur » n'est pas un rôle
stocké mais l'absence de compte ; le « porteur de lien » n'est pas davantage un
rôle : son autorisation tient au seul jeton qu'il présente, ce qui en fait une
frontière de confiance traitée à la section~\ref{sec:security_architecture}.

\section{Méthode de recueil des exigences}
\label{sec:req_method}

Les exigences proviennent de quatre sources :

\begin{enumerate}[leftmargin=*]
  \item le cahier des charges approuvé du projet, qui fixe le périmètre
        initial et les exigences non fonctionnelles chiffrées ;
  \item l'inspection du code livré — modules d'interface, services
        applicatifs, modèles de données, politique d'autorisation — qui établit
        le périmètre réellement implémenté ;
  \item les notes de conception produites pendant la réalisation, qui
        documentent les arbitrages ;
  \item les résultats de la campagne de validation, qui confirment ou
        infirment la satisfaction de chaque exigence.
\end{enumerate}

La priorisation suit un modèle MoSCoW, défini au
tableau~\ref{tab:moscow}.

\begin{table}[htbp]
  \centering
  \caption{Priorités MoSCoW retenues}
  \label{tab:moscow}
  \begin{tabularx}{\textwidth}{L{3.0cm} X}
    \toprule
    \textbf{Priorité} & \textbf{Signification dans ce projet} \\
    \midrule
    \texttt{Doit}    & Indispensable au périmètre central : sans cette capacité, la chaîne de recrutement est interrompue. \\
    \texttt{Devrait} & Important pour la complétude du produit ou pour sa validation, sans bloquer la démonstration minimale. \\
    \texttt{Pourrait}& Capacité optionnelle, réalisable après les exigences principales. \\
    \texttt{Support} & Élément implémenté qui soutient le projet sans constituer une exigence produit mise en avant. \\
    \bottomrule
  \end{tabularx}
\end{table}

\section{Exigences fonctionnelles}
\label{sec:functional_requirements}

Les exigences fonctionnelles sont groupées par domaine. Chaque groupe est
rattaché, au chapitre~\ref{chap:implementation}, à son implémentation, et au
chapitre~\ref{chap:validation}, à son résultat de validation.

\subsection{Comptes, authentification et autorisation}
\label{sec:fr_auth}

\begin{table}[htbp]
  \centering
  \caption{Exigences fonctionnelles — comptes et accès}
  \label{tab:fr_auth}
  \begin{tabularx}{\textwidth}{L{2.4cm} X L{1.9cm}}
    \toprule
    \textbf{ID} & \textbf{Exigence} & \textbf{Priorité} \\
    \midrule
    FR-AUTH-01 & Le système doit permettre l'inscription par adresse et mot de passe, la vérification de l'adresse, la connexion, la réinitialisation et le changement de mot de passe. & \texttt{Doit} \\
    FR-AUTH-02 & Le système doit proposer un second facteur d'authentification fondé sur un mot de passe à usage unique. & \texttt{Doit} \\
    FR-AUTH-03 & Le système doit plafonner les tentatives de connexion et verrouiller temporairement un compte visé. & \texttt{Doit} \\
    FR-AUTH-04 & Le système ne doit pas révéler, par sa réponse ou son temps de réponse, l'existence d'une adresse enregistrée. & \texttt{Doit} \\
    FR-AUTH-05 & Le système devrait permettre la connexion par un fournisseur d'identité externe. & \texttt{Devrait} \\
    FR-ACC-01  & Le système doit gérer quatre rôles — administrateur, responsable RH, recruteur, candidat — avec héritage de permissions. & \texttt{Doit} \\
    FR-ACC-02  & Le système doit refuser par défaut toute action non explicitement autorisée par la politique. & \texttt{Doit} \\
    FR-ACC-03  & Le système doit permettre à un utilisateur de gérer son profil, sa photo et ses coordonnées. & \texttt{Devrait} \\
    FR-ACC-04  & Le système doit journaliser les actions sensibles avec acteur, action et cible. & \texttt{Doit} \\
    FR-ACC-05  & Le système doit permettre à un utilisateur de supprimer son compte. & \texttt{Devrait} \\
    \bottomrule
  \end{tabularx}
\end{table}

Ces exigences fondent le socle. Le refus par défaut (FR-ACC-02) est structurant :
il impose que toute route nouvelle soit explicitement autorisée, plutôt que
d'être ouverte tant qu'on n'a pas pensé à la fermer.

\subsection{Offres d'emploi}
\label{sec:fr_offers}

\begin{table}[htbp]
  \centering
  \caption{Exigences fonctionnelles — offres d'emploi}
  \label{tab:fr_offers}
  \begin{tabularx}{\textwidth}{L{2.4cm} X L{1.9cm}}
    \toprule
    \textbf{ID} & \textbf{Exigence} & \textbf{Priorité} \\
    \midrule
    FR-OFF-01 & Le système doit permettre la création, la modification et la suppression d'une offre d'emploi. & \texttt{Doit} \\
    FR-OFF-02 & Une offre doit porter un intitulé, une description, des compétences requises et un statut de publication. & \texttt{Doit} \\
    FR-OFF-03 & Le système doit exposer publiquement les offres publiées, avec recherche et filtres. & \texttt{Doit} \\
    FR-OFF-04 & Le système doit permettre à un candidat d'enregistrer une offre en favori. & \texttt{Pourrait} \\
    FR-OFF-05 & Le système devrait recommander des offres à un candidat à partir de son profil. & \texttt{Devrait} \\
    \bottomrule
  \end{tabularx}
\end{table}

\subsection{Candidatures et gestion des CV}
\label{sec:fr_applications}

\begin{table}[htbp]
  \centering
  \caption{Exigences fonctionnelles — candidatures et CV}
  \label{tab:fr_applications}
  \begin{tabularx}{\textwidth}{L{2.4cm} X L{1.9cm}}
    \toprule
    \textbf{ID} & \textbf{Exigence} & \textbf{Priorité} \\
    \midrule
    FR-CAND-01 & Le système doit permettre à un candidat de postuler à une offre en joignant un CV. & \texttt{Doit} \\
    FR-CAND-02 & Le système doit contrôler le type et la taille du fichier déposé avant tout traitement. & \texttt{Doit} \\
    FR-CAND-03 & Le système doit permettre l'enregistrement d'un CV de référence réutilisable. & \texttt{Devrait} \\
    FR-CAND-04 & Le système doit exposer au candidat le statut de chacune de ses candidatures. & \texttt{Doit} \\
    FR-CAND-05 & Le système doit permettre à un candidat de retirer une candidature. & \texttt{Devrait} \\
    FR-CAND-06 & Le système doit supprimer le fichier associé lorsqu'une candidature est supprimée. & \texttt{Doit} \\
    FR-CAND-07 & Le système ne doit exposer aucune donnée personnelle de CV dans ses journaux. & \texttt{Doit} \\
    \bottomrule
  \end{tabularx}
\end{table}

L'exigence FR-CAND-07 n'est pas une précaution de confort : un CV concentre
nom, coordonnées, parcours et parfois situation familiale. Sa présence dans un
journal d'exploitation, consultable par des personnes non habilitées,
constituerait une exposition durable et difficile à corriger.

\subsection{Analyse de CV et notation}
\label{sec:fr_scoring}

\begin{table}[htbp]
  \centering
  \caption{Exigences fonctionnelles — analyse de CV et notation}
  \label{tab:fr_scoring}
  \begin{tabularx}{\textwidth}{L{2.4cm} X L{1.9cm}}
    \toprule
    \textbf{ID} & \textbf{Exigence} & \textbf{Priorité} \\
    \midrule
    FR-CV-01    & Le système doit extraire le texte d'un CV déposé en PDF ou en format bureautique. & \texttt{Doit} \\
    FR-CV-02    & Le système doit extraire du texte les compétences, l'expérience et la formation. & \texttt{Doit} \\
    FR-CV-03    & L'analyse doit s'exécuter de façon asynchrone, sans faire attendre le candidat. & \texttt{Doit} \\
    FR-SCORE-01 & Le système doit calculer un score de correspondance entre une candidature et les exigences de l'offre. & \texttt{Doit} \\
    FR-SCORE-02 & Le score doit être décomposé en sous-scores explicables et consultables. & \texttt{Doit} \\
    FR-SCORE-03 & Le système doit classer les candidats d'une offre selon leur score. & \texttt{Doit} \\
    FR-SCORE-04 & Le système doit permettre de contrôler la visibilité du score côté candidat. & \texttt{Devrait} \\
    FR-SCORE-05 & Le système devrait restituer au candidat une lecture compréhensible de l'analyse de son CV. & \texttt{Devrait} \\
    \bottomrule
  \end{tabularx}
\end{table}

FR-SCORE-02 est l'exigence centrale du projet. C'est elle qui distingue PIQBIT
d'un filtre : un score non décomposé ne serait ni justifiable devant un
candidat, ni corrigeable par un recruteur qui le juge erroné.

\subsection{Rapprochement sémantique et assistance}
\label{sec:fr_semantic}

\begin{table}[htbp]
  \centering
  \caption{Exigences fonctionnelles — sémantique et assistant}
  \label{tab:fr_semantic}
  \begin{tabularx}{\textwidth}{L{2.4cm} X L{1.9cm}}
    \toprule
    \textbf{ID} & \textbf{Exigence} & \textbf{Priorité} \\
    \midrule
    FR-SEM-01 & Le système doit indexer les CV analysés sous forme de vecteurs. & \texttt{Doit} \\
    FR-SEM-02 & Le système doit permettre la recherche de profils par proximité de sens. & \texttt{Doit} \\
    FR-SEM-03 & La recherche doit emprunter un index vectoriel natif lorsqu'il est disponible, et un repli fonctionnel sinon. & \texttt{Doit} \\
    FR-SEM-04 & Un échec d'indexation ne doit pas invalider l'analyse ni la notation. & \texttt{Doit} \\
    FR-CAG-01 & Le système doit répondre aux questions internes à partir d'une base de connaissances. & \texttt{Devrait} \\
    FR-CAG-02 & Chaque réponse doit citer les sources dont elle est extraite. & \texttt{Doit} \\
    \bottomrule
  \end{tabularx}
\end{table}

FR-CAG-02 porte une décision de conception assumée : l'assistant est
\emph{extractif} et non génératif. Il restitue des passages issus de la base et
les cite, plutôt que de formuler une réponse dont rien ne garantirait l'ancrage.

\subsection{Entretiens et évaluation}
\label{sec:fr_interviews}

\begin{table}[htbp]
  \centering
  \caption{Exigences fonctionnelles — entretiens et test technique}
  \label{tab:fr_interviews}
  \begin{tabularx}{\textwidth}{L{2.4cm} X L{1.9cm}}
    \toprule
    \textbf{ID} & \textbf{Exigence} & \textbf{Priorité} \\
    \midrule
    FR-INT-01  & Le système doit permettre de proposer jusqu'à cinq créneaux d'entretien à un candidat. & \texttt{Doit} \\
    FR-INT-02  & Le candidat doit pouvoir confirmer un créneau par un lien public, sans compte. & \texttt{Doit} \\
    FR-INT-03  & Le système doit assigner un recruteur à l'entretien et l'avertir de la confirmation. & \texttt{Doit} \\
    FR-INT-04  & Un recruteur ne doit voir que les entretiens qui lui sont affectés. & \texttt{Doit} \\
    FR-INT-05  & Le système doit générer un lien de visioconférence, avec repli si le service principal est indisponible. & \texttt{Doit} \\
    FR-INT-06  & Le système devrait synchroniser l'entretien confirmé avec un agenda externe. & \texttt{Devrait} \\
    FR-INT-07  & Le système doit expirer automatiquement les invitations restées sans réponse. & \texttt{Doit} \\
    FR-INT-08  & Le système doit permettre l'évaluation d'un entretien une fois celui-ci passé. & \texttt{Devrait} \\
    FR-TEST-01 & Le système doit proposer un test technique chronométré par lien public. & \texttt{Doit} \\
    FR-TEST-02 & Le système doit corriger automatiquement le test et en calculer le score. & \texttt{Doit} \\
    FR-TEST-03 & Le seuil d'éligibilité doit être évalué côté serveur. & \texttt{Doit} \\
    FR-TEST-04 & Le système doit clore automatiquement les sessions dont le temps est écoulé. & \texttt{Doit} \\
    \bottomrule
  \end{tabularx}
\end{table}

\subsection{Salaire, négociation et contrats}
\label{sec:fr_contracts}

\begin{table}[htbp]
  \centering
  \caption{Exigences fonctionnelles — salaire, négociation et contrats}
  \label{tab:fr_contracts}
  \begin{tabularx}{\textwidth}{L{2.4cm} X L{1.9cm}}
    \toprule
    \textbf{ID} & \textbf{Exigence} & \textbf{Priorité} \\
    \midrule
    FR-SAL-01  & Le système doit prédire une fourchette salariale à partir d'un modèle entraîné. & \texttt{Doit} \\
    FR-SAL-02  & La prédiction doit s'accompagner d'un intervalle et non d'une valeur unique. & \texttt{Devrait} \\
    FR-NEGO-01 & Le système doit enregistrer les tours successifs d'une négociation avec montants et arguments. & \texttt{Devrait} \\
    FR-NEGO-02 & Le système devrait assister le recruteur dans la formulation d'une contre-proposition. & \texttt{Devrait} \\
    FR-CTR-01  & Le système doit permettre d'éditer un contrat à partir d'une candidature retenue. & \texttt{Doit} \\
    FR-CTR-02  & Le contrat doit être signé par lien public à jeton, sans compte. & \texttt{Doit} \\
    FR-CTR-03  & Le système doit calculer une empreinte des termes signés. & \texttt{Doit} \\
    FR-CTR-04  & Le système doit produire un certificat et une piste d'audit de la signature. & \texttt{Doit} \\
    FR-CTR-05  & Un lien de signature rejoué ne doit jamais signer deux fois. & \texttt{Doit} \\
    FR-CTR-06  & La signature doit créer ou rattacher la fiche employé, sans doublon. & \texttt{Doit} \\
    \bottomrule
  \end{tabularx}
\end{table}

\subsection{Ressources humaines internes et restitution}
\label{sec:fr_hr}

\begin{table}[htbp]
  \centering
  \caption{Exigences fonctionnelles — RH interne, rapports et notifications}
  \label{tab:fr_hr}
  \begin{tabularx}{\textwidth}{L{2.4cm} X L{1.9cm}}
    \toprule
    \textbf{ID} & \textbf{Exigence} & \textbf{Priorité} \\
    \midrule
    FR-RH-01    & Le système doit gérer un répertoire des employés. & \texttt{Devrait} \\
    FR-RH-02    & Le système doit gérer des budgets par département et l'imputation des dépenses. & \texttt{Devrait} \\
    FR-RH-03    & Le système devrait gérer des plans de carrière avec objectifs et jalons. & \texttt{Pourrait} \\
    FR-REP-01   & Le système doit produire des synthèses de l'activité de recrutement. & \texttt{Devrait} \\
    FR-REP-02   & Le système devrait archiver des instantanés de ces synthèses. & \texttt{Pourrait} \\
    FR-REP-03   & Le système devrait proposer une aide à la décision comparant les candidats. & \texttt{Devrait} \\
    FR-NOTIF-01 & Le système doit notifier les acteurs des événements les concernant. & \texttt{Doit} \\
    FR-NOTIF-02 & Les notifications doivent être diffusées en temps réel sur une liaison authentifiée. & \texttt{Devrait} \\
    \bottomrule
  \end{tabularx}
\end{table}

\section{Exigences non fonctionnelles}
\label{sec:nonfunctional_requirements}

Le tableau~\ref{tab:nfr} énonce les exigences non fonctionnelles avec leur
critère d'acceptation et leur état de vérification au
22~septembre~2026. Cet état est repris et démontré au
chapitre~\ref{chap:validation}.

\begin{table}[htbp]
  \centering
  \caption{Exigences non fonctionnelles et état de vérification}
  \label{tab:nfr}
  \begin{tabularx}{\textwidth}{L{1.9cm} L{2.5cm} X L{2.2cm}}
    \toprule
    \textbf{ID} & \textbf{Catégorie} & \textbf{Critère d'acceptation} & \textbf{État} \\
    \midrule
    NFR-01 & Sécurité      & Les mots de passe sont hachés par une fonction lente et salée. & Vérifiée \\
    NFR-02 & Sécurité      & Toute action sensible est refusée par défaut sans autorisation explicite. & \textbf{Écart constaté} \\
    NFR-03 & Sécurité      & Les jetons publics font au moins 32 octets d'entropie. & Vérifiée \\
    NFR-04 & Confidentialité & Aucune donnée personnelle de CV n'apparaît dans les journaux. & Vérifiée \\
    NFR-05 & Fiabilité     & L'indisponibilité d'un service externe ne bloque aucun parcours. & Vérifiée \\
    NFR-06 & Intégrité     & Le schéma de la base est piloté exclusivement par migrations. & Vérifiée \\
    NFR-07 & Performance   & Le dépôt d'une candidature répond sans attendre l'analyse du CV. & Vérifiée \\
    NFR-08 & Performance   & Le bundle initial du client reste sous le budget fixé. & \textbf{Non satisfaite} \\
    NFR-09 & Maintenabilité & Le code est typé et validé aux frontières. & Vérifiée \\
    NFR-10 & Disponibilité & Le service atteint 99,5\,\% de disponibilité en production. & \textbf{Non vérifiée} \\
    NFR-11 & Accessibilité & L'interface respecte les critères d'accessibilité retenus. & \textbf{Non vérifiée} \\
    \bottomrule
  \end{tabularx}
\end{table}

Trois exigences ne sont pas satisfaites, et le rapport les traite comme telles.
NFR-08 est mesurée en dépassement. NFR-10 et NFR-11 ne sont pas vérifiées faute
d'environnement de production et de campagne d'accessibilité : les déclarer
vérifiées serait une affirmation sans preuve. NFR-02 présente un écart révélé
par la campagne de validation, analysé à la
section~\ref{sec:failure_analysis}.

\section{Cas d'utilisation}
\label{sec:use_cases}

La figure~\ref{fig:use_cases} donne la vue d'ensemble des cas d'utilisation et
de leurs acteurs.

\pqfigure{use_cases}
  {Diagramme de cas d'utilisation de PIQBIT : acteurs, capacités et héritage
   de permissions entre rôles internes.}
  {fig:use_cases}

\subsection{Décomposition détaillée des cas principaux}
\label{sec:use_cases_detail}

Le tableau~\ref{tab:use_case_catalog} catalogue les cas d'utilisation
principaux, puis quatre d'entre eux sont détaillés.

\begin{table}[htbp]
  \centering
  \caption{Catalogue des cas d'utilisation principaux}
  \label{tab:use_case_catalog}
  \begin{tabularx}{\textwidth}{L{1.7cm} L{4.6cm} L{3.0cm} X}
    \toprule
    \textbf{ID} & \textbf{Cas d'utilisation} & \textbf{Acteur principal} & \textbf{Exigences} \\
    \midrule
    UC-01 & Postuler à une offre                  & Candidat        & FR-CAND-01/02 \\
    UC-02 & Analyser et noter une candidature     & Traitement de fond & FR-CV-01/02/03, FR-SCORE-01/02/03 \\
    UC-03 & Consulter le classement des candidats & Recruteur       & FR-SCORE-03 \\
    UC-04 & Rechercher un profil par le sens      & Recruteur       & FR-SEM-01/02/03 \\
    UC-05 & Planifier un entretien                & Recruteur       & FR-INT-01/03 \\
    UC-06 & Confirmer un créneau d'entretien      & Porteur de lien & FR-INT-02/05 \\
    UC-07 & Passer le test technique              & Porteur de lien & FR-TEST-01/02 \\
    UC-08 & Estimer un salaire et négocier        & Responsable RH  & FR-SAL-01, FR-NEGO-01 \\
    UC-09 & Signer un contrat                     & Porteur de lien & FR-CTR-02/03/04/05 \\
    UC-10 & Activer la fiche employé              & Système         & FR-CTR-06 \\
    UC-11 & Administrer comptes et rôles          & Administrateur  & FR-ACC-01/02/04 \\
    \bottomrule
  \end{tabularx}
\end{table}

\paragraph{UC-01 — Postuler à une offre.}
\emph{Acteur :} candidat authentifié. \emph{Précondition :} l'offre est
publiée et le candidat n'y a pas déjà postulé. \emph{Scénario nominal :} le
candidat ouvre la fiche de l'offre, joint un CV ou réutilise son CV de
référence, valide ; le système contrôle le type et la taille du fichier,
enregistre la candidature, planifie l'analyse et répond immédiatement.
\emph{Variantes :} format non accepté ou taille excessive — la candidature est
refusée avec un message explicite. \emph{Postcondition :} la candidature existe
au statut initial et son analyse est planifiée.

\paragraph{UC-02 — Analyser et noter une candidature.}
\emph{Acteur :} traitement de fond. \emph{Précondition :} une candidature vient
d'être déposée. \emph{Scénario nominal :} le traitement extrait le texte du CV,
en tire compétences, expérience et formation, enregistre l'analyse, indexe les
vecteurs, calcule le score et reclasse les candidats de l'offre.
\emph{Variantes :} l'indexation échoue — la notation est conservée, la
recherche sémantique sera simplement moins riche. \emph{Postcondition :} la
candidature porte une analyse, un score décomposé et un rang.

\paragraph{UC-06 — Confirmer un créneau d'entretien.}
\emph{Acteur :} porteur de lien, sans compte. \emph{Précondition :} une
invitation est active et non expirée. \emph{Scénario nominal :} le porteur
ouvre le lien reçu, consulte les créneaux, en choisit un ; le système vérifie
que le créneau n'a pas été réservé entre-temps, marque l'invitation confirmée,
avertit le recruteur assigné et synchronise l'agenda.
\emph{Variantes :} lien expiré, déjà confirmé ou annulé — l'action est refusée
avec une raison distincte dans chaque cas. \emph{Postcondition :} l'entretien
est confirmé et le recruteur assigné en est informé.

\paragraph{UC-09 — Signer un contrat.}
\emph{Acteur :} porteur de lien, sans compte. \emph{Précondition :} un contrat
a été envoyé et n'est ni signé, ni refusé. \emph{Scénario nominal :} le porteur
ouvre le lien, consulte les termes en lecture seule, saisit son nom et son
tracé de signature ; le système vérifie l'état du contrat, calcule l'empreinte
des termes, émet un certificat, enregistre la piste d'audit, puis crée ou
rattache la fiche employé. \emph{Variantes :} contrat déjà signé — refus
explicite ; lien rejoué — aucune seconde signature. \emph{Postcondition :} le
contrat est signé, certifié, et la fiche employé est active.

\section{Acceptation et traçabilité}
\label{sec:traceability}

Chaque exigence est reliée à son implémentation et à sa validation selon le
modèle du tableau~\ref{tab:traceability_model}. La matrice complète figure en
annexe~\ref{app:traceability}.

\begin{table}[htbp]
  \centering
  \caption{Modèle de traçabilité des exigences}
  \label{tab:traceability_model}
  \begin{tabularx}{\textwidth}{L{3.4cm} X}
    \toprule
    \textbf{Élément} & \textbf{Contenu} \\
    \midrule
    Identifiant d'exigence & Référence stable, réutilisée sans modification aux chapitres 4 et 5 \\
    Domaine                & Groupe fonctionnel d'appartenance \\
    Implémentation         & Sous-système du chapitre~\ref{chap:implementation} qui réalise l'exigence \\
    Validation             & Résultat daté du chapitre~\ref{chap:validation} \\
    État                   & Satisfaite, partiellement satisfaite, non satisfaite ou non vérifiée \\
    \bottomrule
  \end{tabularx}
\end{table}

Le vocabulaire d'état est volontairement restreint. « Non vérifiée » n'est pas
un synonyme atténué de « satisfaite » : il signale qu'aucune preuve n'existe,
ce qui est une information en soi.

\section{Limites de périmètre}
\label{sec:scope_boundaries}

Les éléments suivants sont explicitement hors périmètre :

\begin{itemize}[leftmargin=*]
  \item l'analyse de CV en langues autres que celles couvertes par le modèle
        de plongement retenu ;
  \item la détection de fraude documentaire sur les CV ;
  \item l'intégration avec des jobboards externes pour la diffusion des offres ;
  \item la paie et la gestion administrative des employés au-delà de la fiche
        et du budget ;
  \item la conformité juridique de la signature électronique au sens d'un
        règlement de signature qualifiée : la signature implémentée est
        auto-hébergée, avec empreinte et piste d'audit, ce qui en fait une
        preuve d'intention et non une signature qualifiée.
\end{itemize}

\section{Conclusion du chapitre}
\label{sec:req_conclusion}

Ce chapitre a identifié les acteurs de PIQBIT, formalisé les exigences
fonctionnelles en huit domaines et les exigences non fonctionnelles avec leurs
critères d'acceptation, décrit les cas d'utilisation principaux et posé le
modèle de traçabilité. Trois exigences non fonctionnelles sont d'ores et déjà
identifiées comme non satisfaites ou non vérifiées, ce que le
chapitre~\ref{chap:validation} documentera par des mesures datées.

Le chapitre suivant traduit ces exigences en architecture : contexte système,
découpage en composants, modèle de données, séquences principales, frontières
de sécurité et décisions d'architecture justifiées.
